\subsection{Penance and Mercy}\label{sec:penance}

The Code of Canon Law of 1983, which governs the
Latin Church throughout the world, teaches that all the faithful are bound by
divine law to do penance, each in his own way, and it names the penitential
days and times of the whole Church: every Friday of the year, and the season of
Lent (cc.~1249--1250). Advent is not among them. The Code prescribes
abstinence on Fridays and both fast and abstinence on Ash Wednesday and Good
Friday, and it allows each conference of bishops to set the observance more
precisely or to put other works of penance, above all works of charity and
exercises of piety, in place of abstinence and fast (cc.~1251--1253).

In the United States the bishops issued their Pastoral Statement on Penance and
Abstinence on 18 November 1966, to take effect on the First Sunday of Advent of
that year. It asks Catholics to make Advent a time of meditation on what the
liturgy teaches and of fuller participation in the rites in which the mysteries
of Advent are celebrated (n.~7). It imposes no fast before any feast, and it
states that vigils and Ember days no longer bind to fast and abstinence
(n.~17). In 1983 the bishops declared that these norms remain in force under
the new Code. The Universal Norms on the Liturgical Year left the Ember days to
be arranged by each conference of bishops (n.~46), and the present calendar of
the United States has none. So in this country no fast or abstinence is
required for Advent as such.

The fasts of earlier ages, the three weekly fasts of Tours and Mâcon from St
Martin's day to Christmas, belong to history, and none of them is law now. The
same 1966 statement leaves the devout free to prepare for a feast by a day of self-denial, prayer and fasting. A Christian may well choose
to keep some fast in Advent, privately and gladly. The question is how to keep
it, and the Fathers are clear that a fast is not kept by the stomach alone.

Pope St Leo the Great, preaching to Rome at the Ember fast of December
(Section~\ref{sec:history}), told his people what any fast is for:

\begin{quote}
When the body therefore fasts from food, let the mind fast from vices, and pass
judgment upon all earthly cares and desires according to the law of its King\par
{\smallskip\noindent\small ---St Leo the Great, Sermon 19\par}
\end{quote}

\noindent And he told them where the food saved should go:

\begin{quote}
And hence we warn you, beloved, in fatherly affection, to make this winter fast
fruitful to yourselves by bounteous alms, rejoicing that by you the Lord feeds
and clothes His poor, to whom assuredly He could have given the possessions
which He has bestowed on you, had He not in His unspeakable mercy wished to
justify them for their patient labour, and you for your works of love.\par
{\smallskip\noindent\small ---St Leo the Great, Sermon 19\par}
\end{quote}

The Lord himself has said what he will ask at his coming:

\begin{quote}
For I was hungry, and you gave me to eat: I was thirsty, and you gave me to
drink: I was a stranger, and you took me in: Naked, and you covered me: sick,
and you visited me: I was in prison, and you came to me. \ldots{} Amen I say to
you, as long as you did it to one of these my least brethren, you did it to
me.\par
{\smallskip\noindent\small ---Matthew 25:35--36, 40\par}
\end{quote}

\noindent St Cyril of Jerusalem, in the same lecture that the Church reads on
the First Sunday, came to this page and would not let his hearers turn it into
a figure: ``Learn hence the way; there is here no need of allegory, but to
fulfil what is said.'' A little later he added: ``At once then begin to do these
works, and abide in the faith; lest, like the foolish virgins, tarrying to buy
oil, thou be shut out.''

St John Chrysostom asked why the King calls the poor his least brethren, and
whom he means:

\begin{quote}
What sayest Thou? they are Thy brethren; and how dost Thou call them least. Why,
for this reason they are brethren, because they are lowly, because they are
poor, because they are outcast. For such doth He most invite to brotherhood, the
unknown, the contemptible, not meaning by these the monks only, and them that
have occupied the mountains, but every believer; though he be a secular person,
yet if he be hungry, and famishing, and naked, and a stranger, His will is he
should have the benefit of all this care. For baptism renders a man a brother,
and the partaking of the divine mysteries.\par
{\smallskip\noindent\small ---St John Chrysostom, Homily 79 on St Matthew\par}
\end{quote}

There is also a mercy that is borne rather than given. On
the Saturday of the first week the Office reads St Cyprian, bishop of Carthage
and martyr, on patience, and he binds patience and charity together:

\begin{quote}
Charity is the bond of brotherhood, the foundation of peace, the holdfast and
security of unity, which is greater than both hope and faith, which excels both
good works and martyrdoms, which will abide with us always, eternal with God in
the kingdom of heaven. Take from it patience; and deprived of it, it does not
endure. Take from it the substance of bearing and of enduring, and it continues
with no roots nor strength.\par
{\smallskip\noindent\small ---St Cyprian, \emph{On the Advantage of Patience}\par}
\end{quote}

\noindent The December weeks are crowded and short-tempered, and the patience
Cyprian asks for is a penance no law imposes and every household needs. Give
alms in Advent from what the season tempts you to spend on yourself; visit
someone who will be alone at Christmas; bear the delays and the crowds without
complaint.
